\chapter{Concurrency, networking, and distributed systems}

\section{Concurrency is about interleaving}
Concurrency means multiple activities can make progress during overlapping
intervals. Parallelism is simultaneous execution. A single-threaded event loop
can be concurrent without being parallel; multiple CPU cores can be parallel.
The distinction matters for reasoning, resource limits, and performance claims.

The basic danger is an interleaving that violates an invariant. Suppose two
requests both read an account balance of 100 and both decide that a withdrawal
of 80 is allowed. If each writes 20, the final balance looks plausible but the
system has accepted 160 of withdrawals. A transaction, an atomic conditional
update, or a serialized command can protect the rule.

\begin{keyidea}
Do not ask only whether each operation is correct in isolation. Ask which
interleavings are possible, which state is shared, and which invariant must hold
after every visible transition.
\end{keyidea}

\section{Memory and message boundaries}
In one process, threads share memory and require synchronization. Across
processes, machines, or organizations, messages cross a boundary and can be
delayed, duplicated, reordered, or lost. A network timeout does not tell the
caller whether the remote operation failed or succeeded but its response was
lost. This is the source of many duplicate writes.

A robust request path defines:

\begin{itemize}
  \item a deadline propagated through every downstream call;
  \item a bounded retry policy only for operations that are safe to repeat;
  \item an idempotency key or deduplication record for retried effects;
  \item a clear failure result when the outcome is unknown;
  \item telemetry that distinguishes timeout, refusal, overload, and validation.
\end{itemize}

Retries multiply load during failure. Exponential backoff with jitter spreads
attempts, but it is not a cure for an overloaded dependency. A caller should
retry only when the error is plausibly transient and the total deadline still
allows useful work. A circuit breaker can stop calls to a failing dependency,
but it introduces its own state machine and false-open trade-off.

\section{Ordering and causality}
Wall-clock timestamps are not a reliable global ordering in a distributed system.
Clocks differ and can move. Lamport's work on logical clocks shows how a system
can represent a partial order of events without pretending that every event has
one exact global time \cite{lamport1978time}. In practice, choose the smallest
ordering guarantee that the product needs: per-account sequence numbers,
partition-local order, or an explicit version check.

\section{Consistency is a product decision}
An eventually consistent view can be correct for a feed whose purpose is
discovery, but unacceptable for a bank balance or an authorization decision.
“Strong” and “eventual” are not complete specifications. State what a reader may
observe after a write, which conflicts can occur, and how a user can repair or
understand them.

The CAP theorem is frequently reduced to a slogan about choosing two of three
properties. The more useful interpretation is narrower: under a network
partition, a distributed data system cannot simultaneously provide a particular
strong consistency guarantee and uninterrupted availability for every operation.
The actual design space includes latency, scope of consistency, failure model,
and recovery strategy \cite{brewer2012cap}.

\section{Queues and workers}
Queues decouple producers and consumers in time. They absorb bursts and allow a
worker pool to process work at a controlled rate. They also create backlog,
duplicate delivery, poison messages, and delayed user-visible state. A queue is
not a free buffer; it is a new stateful subsystem.

For a message handler, specify:

\begin{enumerate}
  \item what makes a message valid;
  \item the acknowledgment point;
  \item the effect of a crash before and after acknowledgment;
  \item maximum attempts and dead-letter handling;
  \item deduplication identity and retention;
  \item how operators replay, quarantine, or delete a message safely.
\end{enumerate}

\section{A small state machine}
Model a payment or order lifecycle as a state machine before implementing it:

\begin{center}
\begin{tikzpicture}[node distance=20mm,>=Latex]
  \node[draw,rounded corners,fill=codebg,minimum width=25mm,minimum height=9mm] (new) {new};
  \node[draw,rounded corners,fill=codebg,minimum width=25mm,minimum height=9mm,right=of new] (accepted) {accepted};
  \node[draw,rounded corners,fill=codebg,minimum width=25mm,minimum height=9mm,right=of accepted] (paid) {paid};
  \node[draw,rounded corners,fill=codebg,minimum width=25mm,minimum height=9mm,below=of accepted] (cancelled) {cancelled};
  \draw[->,thick,accent] (new) -- node[above,font=\scriptsize]{validate} (accepted);
  \draw[->,thick,accent] (accepted) -- node[above,font=\scriptsize]{capture} (paid);
  \draw[->,thick,accent] (accepted) -- node[right,font=\scriptsize]{cancel} (cancelled);
\end{tikzpicture}
\end{center}

The diagram makes illegal transitions visible. It also raises operational
questions: what if capture succeeds but the response times out? The answer may
be a reconciliation state rather than an optimistic retry.

\begin{exercise}
For a queue consumer that sends an email, list the states before and after a
crash at each of these points: before send, after send, before acknowledgment,
and after acknowledgment. Choose a deduplication strategy and state its limits.
\end{exercise}

\section{Chapter summary}
Concurrency is reasoning about interleavings; distribution adds uncertain
messages and partial failure. Define deadlines, retryability, idempotency,
ordering, consistency, and recovery. Treat queues and state machines as explicit
parts of the product rather than invisible implementation details.

